\chapter{Correspondência com o plano curricular}
\label{ap:curriculo}

Este apêndice estabelece a correspondência entre as unidades curriculares do mestrado e os pontos do
trabalho em que a respetiva matéria foi aplicada. A Tabela~\ref{tab:ap_curriculo} apresenta essa
correspondência e identifica igualmente as unidades cuja matéria não teve aplicação, uma vez que uma
correspondência que encontre aplicação para todas as unidades em todos os pontos não transmite
informação.

Duas entradas da tabela declaram ausência de aplicação: a matéria de Ambientes Inteligentes e a
componente de programação lógica e por restrições. O seu registo é mais informativo do que a
construção de uma ligação artificial.

\begin{table}[!htbp]
\caption[Correspondência com o plano curricular]{Aplicação da matéria de cada unidade curricular do
plano de estudos do mestrado no trabalho descrito nesta dissertação. A última coluna indica a secção
em que essa aplicação é visível.}
\label{tab:ap_curriculo}
\centering
\footnotesize
\begin{tabular}{L{0.23\textwidth} L{0.47\textwidth} L{0.19\textwidth}}
\toprule
\tabhead{Unidade} & \tabhead{Onde é aplicada} & \tabhead{Secção} \\
\midrule
\multicolumn{3}{l}{\emph{Fundamentos de sistemas inteligentes}} \\
Aprendizagem Automática 1 &
  Construção de um conjunto de treino rotulado a partir de duas fontes desalinhadas no tempo;
  regressão logística com normalização e reponderação de classes; conjuntos de árvores por reforço
  do gradiente como referência superior; calibração \emph{a posteriori} &
  §\ref{sec:met_triagem}, §\ref{sec:av_qi3} \\
Engenharia do Conhecimento &
  Raciocínio baseado em casos: recuperação e reutilização de precedentes em substituição de uma
  regra geral, com a interrupção deliberada do ciclo antes da adaptação &
  §\ref{sec:ctx_casos}, §\ref{sec:met_recuperacao} \\
Paradigmas de Programação em \gls{IA} &
  Separação entre lógica pura e efeitos externos, que é a propriedade que permite verificar o núcleo
  do sistema sem acesso à rede e sem carregar modelos. A matéria de programação lógica e por
  restrições não teve aplicação &
  §\ref{sec:met_metodologia} \\
Planeamento e Tomada de Decisão &
  Política de decisão sob orçamento: portas em cascata, limiar derivado de uma análise de custos de
  erro, e distinção entre uma política que dispõe do dia completo e uma política que decide a cada
  ciclo &
  §\ref{sec:sis_politica}, §\ref{sec:av_qi3} \\
\midrule
\multicolumn{3}{l}{\emph{Inteligência artificial avançada e aplicações especializadas}} \\
Aprendizagem Automática 2 &
  Representações contextuais de frase produzidas por uma arquitetura \emph{Transformer}; exportação
  do codificador para um formato leve, com medição da fidelidade da conversão, imposta pela
  restrição de memória do ambiente de execução &
  §\ref{sec:ctx_texto}, §\ref{sec:met_recuperacao} \\
Linguagem Natural e Sistemas Conversacionais &
  Semelhança semântica entre frases e avaliação de recuperação sensível à ordem. O sistema constitui
  uma arquitetura de recuperação sem geração: recupera evidência e recusa produzir texto sobre ela,
  com a justificação apresentada no estado da arte &
  §\ref{sec:ctx_texto}, §\ref{sec:av_qi2} \\
Ambientes Inteligentes &
  Sem aplicação. O trabalho não comporta componente física, robótica nem de Internet das Coisas &
  --- \\
\midrule
\multicolumn{3}{l}{\emph{Profissionalismo, segurança e inovação}} \\
Aspetos Sociais e Éticos em \gls{IA} &
  Confiança apropriada em sistemas automáticos e risco de a explicação aumentar a aceitação de
  respostas incorretas; recusa de aconselhamento; fadiga de alertas como requisito de desenho &
  §\ref{sec:ctx_xai}, §\ref{sec:met_desafios} \\
Privacidade e Segurança em Sistemas Inteligentes &
  Minimização de dados por desenho, com a exceção declarada; gestão de credenciais e mascaramento de
  segredos em registos; verificação de integridade do modelo transferido &
  §\ref{sec:met_privacidade}, §\ref{sec:met_seguranca} \\
Investigação e Inovação em \gls{IA} &
  Enquadramento em investigação por desenho; critérios de sucesso fixados antes das medições; linhas
  de comparação declaradas com o respetivo valor de referência; matriz de evidência com as
  afirmações retiradas &
  §\ref{sec:met_metodologia}, Apêndice~\ref{ap:reprodutibilidade} \\
\bottomrule
\end{tabular}
\end{table}
